\section{Problems and solutions}

\begin{problem}[12 points]{A numbered problem}
  Problem statements may contain prose, displayed mathematics, lists, and references.
  The optional argument records the available points.
  \begin{subproblems}
    \item The first subproblem.
    \item The second subproblem.
  \end{subproblems}

  \begin{solution}
    Solutions use a distinct breakable panel and may span multiple pages.
  \end{solution}

  \begin{answer}
    Use an answer box when the final result should be easy to locate.
  \end{answer}
\end{problem}

\begin{problem}{A problem without points}
  The points argument is optional, but the descriptive title is required.
\end{problem}

\subsection*{Lists}

Bullets and numbers both take the accent on the marker and leave the text in ink, so a list reads as structure without the words changing colour.

\begin{itemize}
  \item A first point, at the top level.
  \item A second one, with a nested list under it.
  \begin{itemize}
    \item The second level takes a dash.
    \item Only the marker is coloured, at every depth.
    \begin{itemize}
      \item And the third an asterisk.
    \end{itemize}
  \end{itemize}
  \item Back at the top level.
\end{itemize}

\begin{enumerate}
  \item A numbered list matches the lettered \code{subproblems} above.
  \item Both are set once, in \filepath{theme/ricardo-homework.sty}.
\end{enumerate}
